% GENERATED FILE. Edit canonical lessons, then run npm run generate:books.
% canonical-source-hash: fnv1a64:30b44e1d4bafbfa0
% canonical-lessons: TA-C207-appatiyanal, TA-C207-kittattatta, TA-C207-unmaiyil, TA-C207-potuvaka, TA-C207-valakkamaka

\chapter{In that case, Almost, Actually, In general, Usually}
\label{ch:ta-in-that-case-almost-actually-in-general-usually}
\hlchaptermodality{207}

\begin{chapteropening}
\textbf{By the end of this chapter:} \emph{I can say in that case, almost, actually, in general and usually.}

Complete the last lesson of chapter 207: I can say in that case, almost, actually, in general and usually.
\end{chapteropening}

\section[\texorpdfstring{\ta{அப்படியானால்} (appaṭiyāṉāl)}{அப்படியானால் (appaṭiyāṉāl)}]{\ta{அப்படியானால்} (appaṭiyāṉāl) — in that case}

\label{lesson:TA-C207-appatiyanal}



\begin{quote}
Before the new one: say the Tamil for a playground, then the Tamil for a swimming pool.
\end{quote}

\subsection*{You'll want to know: \ta{அப்படியானால்}}
\textbf{\ta{அப்படியானால்}} — \emph{appaṭiyāṉāl} — \textquotedblleft{}in that case\textquotedblright{}.

\begin{grammarlens}[title={What you've built}]
One more word for this chapter.
\end{grammarlens}

\subsection*{Your turn}
\begin{itemize}
\raggedright
  \item \emph{Say it:} \emph{appaṭiyāṉāl}
  \item \emph{Say it:} \emph{appaṭiyāṉāl}, once more
  \item \emph{Say it:} say \emph{appaṭiyāṉāl}
  \item \emph{Recall:} say the Tamil for a playground, then the Tamil for a swimming pool, then say \emph{appaṭiyāṉāl} again
\end{itemize}

\begin{tcolorbox}[breakable,colback=teal!4,colframe=teal!35!black,title={Before you move on}]
What does \ta{அப்படியானால்} mean? (\textquotedblleft{}In that case\textquotedblright{}.) Say it once more.
\end{tcolorbox}
\section[\texorpdfstring{\ta{கிட்டத்தட்ட} (kiṭṭattaṭṭa)}{கிட்டத்தட்ட (kiṭṭattaṭṭa)}]{\ta{கிட்டத்தட்ட} (kiṭṭattaṭṭa) — almost, nearly}

\label{lesson:TA-C207-kittattatta}



\begin{quote}
Before the new one: say the Tamil for a swimming pool, then the Tamil for in that case.
\end{quote}

\subsection*{You'll want to know: \ta{கிட்டத்தட்ட}}
\textbf{\ta{கிட்டத்தட்ட}} — \emph{kiṭṭattaṭṭa} — \textquotedblleft{}almost, nearly\textquotedblright{}.

\begin{grammarlens}[title={What you've built}]
One more word for this chapter.
\end{grammarlens}

\subsection*{Your turn}
\begin{itemize}
\raggedright
  \item \emph{Say it:} \emph{kiṭṭattaṭṭa}
  \item \emph{Say it:} \emph{kiṭṭattaṭṭa}, once more
  \item \emph{Say it:} say \emph{kiṭṭattaṭṭa}
  \item \emph{Recall:} say the Tamil for a swimming pool, then the Tamil for in that case, then say \emph{kiṭṭattaṭṭa} again
\end{itemize}

\begin{tcolorbox}[breakable,colback=teal!4,colframe=teal!35!black,title={Before you move on}]
What does \ta{கிட்டத்தட்ட} mean? (\textquotedblleft{}Almost, nearly\textquotedblright{}.) Say it once more.
\end{tcolorbox}
\section[\texorpdfstring{\ta{உண்மையில்} (uṇmaiyil)}{உண்மையில் (uṇmaiyil)}]{\ta{உண்மையில்} (uṇmaiyil) — actually, in fact}

\label{lesson:TA-C207-unmaiyil}



\begin{quote}
Before the new one: say the Tamil for in that case, then the Tamil for almost.
\end{quote}

\subsection*{You'll want to know: \ta{உண்மையில்}}
\textbf{\ta{உண்மையில்}} — \emph{uṇmaiyil} — \textquotedblleft{}actually, in fact\textquotedblright{}.

\begin{grammarlens}[title={What you've built}]
One more word for this chapter.
\end{grammarlens}

\subsection*{Your turn}
\begin{itemize}
\raggedright
  \item \emph{Say it:} \emph{uṇmaiyil}
  \item \emph{Say it:} \emph{uṇmaiyil}, once more
  \item \emph{Say it:} say \emph{uṇmaiyil}
  \item \emph{Recall:} say the Tamil for in that case, then the Tamil for almost, then say \emph{uṇmaiyil} again
\end{itemize}

\begin{tcolorbox}[breakable,colback=teal!4,colframe=teal!35!black,title={Before you move on}]
What does \ta{உண்மையில்} mean? (\textquotedblleft{}Actually, in fact\textquotedblright{}.) Say it once more.
\end{tcolorbox}
\section[\texorpdfstring{\ta{பொதுவாக} (poṭuvāka)}{பொதுவாக (poṭuvāka)}]{\ta{பொதுவாக} (poṭuvāka) — in general, generally}

\label{lesson:TA-C207-potuvaka}



\begin{quote}
Before the new one: say the Tamil for almost, then the Tamil for actually.
\end{quote}

\subsection*{You'll want to know: \ta{பொதுவாக}}
\textbf{\ta{பொதுவாக}} — \emph{poṭuvāka} — \textquotedblleft{}in general, generally\textquotedblright{}.

\begin{grammarlens}[title={What you've built}]
One more word for this chapter.
\end{grammarlens}

\subsection*{Your turn}
\begin{itemize}
\raggedright
  \item \emph{Say it:} \emph{poṭuvāka}
  \item \emph{Say it:} \emph{poṭuvāka}, once more
  \item \emph{Say it:} say \emph{poṭuvāka}
  \item \emph{Recall:} say the Tamil for almost, then the Tamil for actually, then say \emph{poṭuvāka} again
\end{itemize}

\begin{tcolorbox}[breakable,colback=teal!4,colframe=teal!35!black,title={Before you move on}]
What does \ta{பொதுவாக} mean? (\textquotedblleft{}In general, generally\textquotedblright{}.) Say it once more.
\end{tcolorbox}
\section[\texorpdfstring{\ta{வழக்கமாக} (vaḻakkamāka)}{வழக்கமாக (vaḻakkamāka)}]{\ta{வழக்கமாக} (vaḻakkamāka) — usually, as usual}

\label{lesson:TA-C207-valakkamaka}



\begin{quote}
Before the new one: say the Tamil for actually, then the Tamil for in general.
\end{quote}

\subsection*{You'll want to know: \ta{வழக்கமாக}}
\textbf{\ta{வழக்கமாக}} — \emph{vaḻakkamāka} — \textquotedblleft{}usually, as usual\textquotedblright{}.

\begin{grammarlens}[title={What you've built}]
One more word for this chapter.
\end{grammarlens}

\subsection*{Your turn}
\begin{itemize}
\raggedright
  \item \emph{Say it:} \emph{vaḻakkamāka}
  \item \emph{Say it:} \emph{vaḻakkamāka}, once more
  \item \emph{Say it:} say \emph{vaḻakkamāka}
  \item \emph{Recall:} say the Tamil for actually, then the Tamil for in general, then say \emph{vaḻakkamāka} again
\end{itemize}

\begin{tcolorbox}[breakable,colback=teal!4,colframe=teal!35!black,title={Before you move on}]
What does \ta{வழக்கமாக} mean? (\textquotedblleft{}Usually, as usual\textquotedblright{}.) Say it once more.
\end{tcolorbox}
